\section{Background}
\subsection{Protein fitness landscapes}
Proteins fold into 3D structures; function emerges from structure
Multiple objectives: stability, yield, activity, thermotolerance

- Most often multiple objectives (folding, temperature, yield)
- Cannot just learn one physical property and generalize
- fitness landscape is the map from sequence space to fitness value
   - epistasis — mutations interact non-additively (add note here + saddle)
   - ruggedness — many local optima
   - sparsity — most sequences are non-functional
\subsection{Search space}
- Ok - so where are we playing. Field is intractable so need to make it more sparse
- Wet lab vs non wet lab
- Full sequence space is intractable 
- Either look at few mutations away max, or look at combinatorial site landscapes
\subsection{Search methods}
--Turbo OG is just breeding
- **Classical directed evolution:** mutate → express → assay → select
      ; repeat
   43   - works but greedy, gets stuck at local optima, expensive per itera
      tion
   44 - **MLDE:** train surrogate on labeled variants → predict fitness acr
      oss unseen space → select next batch
   45   - reduces experimental burden by screening computationally first
   46   - one or multiple rounds
   47 - **ALDE:** explicitly iterative MLDE with uncertainty-aware acquisit
      ion
   48   - acquisition function decides what to measure next (EI, PI, Thomps
      on sampling, UCB)
   49   - active learning loop: label → train → acquire → measure → repeat
   50 - Sequence representations: sequences must be converted to numerical
      features before any ML model can use them — the choice of representat
      ion is a core design decision (elaborated in Methods)


\subsection{The ALDE loop}

Active learning directed evolution (ALDE) is an iterative optimisation strategy that interleaves computational modelling with wet-lab measurement. The core idea is to spend the experimental budget where the model believes it will be most informative or most rewarding, rather than screening sequences blindly.

Every ALDE method is built from four components: an \textbf{initialisation strategy} that selects the first batch of variants to measure before any model is available; a \textbf{sequence encoding} that converts amino acid sequences into numerical feature vectors a machine learning model can consume; a \textbf{surrogate model} trained on the labeled observations to predict the fitness of unmeasured variants; and an \textbf{acquisition function} that uses the surrogate's predictions — and, where available, its uncertainty estimates — to select which variants to measure next.

Given these components, the algorithm proceeds as follows:

\begin{enumerate}
  \item \textbf{Initialise.} Select a first batch of sequences using the initialisation strategy and measure their fitness in the lab.
  \item \textbf{Encode.} Convert all sequences in the candidate pool to feature vectors using the chosen encoding.
  \item \textbf{Train.} Fit the surrogate model on the labeled observations accumulated so far.
  \item \textbf{Acquire.} Score all unlabeled candidates with the acquisition function and select the top batch to measure next.
  \item \textbf{Measure.} Send the selected batch to the lab, obtain fitness values, and add them to the labeled set.
  \item \textbf{Repeat} from step 3 until the experimental budget is exhausted.
\end{enumerate}

The choice of which components to use — and how they interact — determines how efficiently the budget is spent. Section~\ref{sec:methods} gives the formal treatment and details the specific components evaluated in this thesis.
